% !TEX program = xelatex
% 공통수학2 · Ⅲ. 함수와 그래프 · 1. 함수 · 01 함수 · 1/2차시
\documentclass[11pt,a4paper]{article}
\usepackage{kotex}
\usepackage{iftex}
\ifPDFTeX\else
  \setmainhangulfont{Noto Serif CJK KR}
  \setsanshangulfont{Noto Sans CJK KR}
\fi
\usepackage[margin=2.1cm]{geometry}
\usepackage{parskip}
\usepackage{enumitem}
\usepackage{array}
\usepackage{tabularx}
\usepackage{booktabs}
\usepackage{xcolor}
\usepackage{amssymb}
\usepackage{tikz}
\usepackage[most]{tcolorbox}
\usepackage{fancyhdr}
\usepackage{titlesec}

\definecolor{goldc}{HTML}{B07C22}
\definecolor{goldstrong}{HTML}{8A5F16}
\definecolor{indigoc}{HTML}{2B3B57}
\definecolor{indigosoft}{HTML}{6E82A6}
\definecolor{linec}{HTML}{D6DACB}
\definecolor{surface2}{HTML}{E4E8DC}
\definecolor{notebg}{HTML}{FBF3E3}
\definecolor{inksoft}{HTML}{52604F}

\titleformat{\section}{\normalfont\sffamily\bfseries\large\color{indigoc}}{\thesection}{0.6em}{}
\titlespacing{\section}{0pt}{16pt}{8pt}
\newtcolorbox{ideabox}{enhanced, breakable, colback=white, colframe=goldc, boxrule=0.5pt, leftrule=3pt, arc=2pt, left=10pt, right=10pt, top=8pt, bottom=8pt}
\newtcolorbox{calloutbox}{enhanced, breakable, colback=white, colframe=indigosoft, boxrule=0.5pt, leftrule=3pt, arc=2pt, left=10pt, right=10pt, top=7pt, bottom=7pt}
\newtcolorbox{notebox}{enhanced, breakable, colback=notebg, colframe=goldc, boxrule=0.5pt, leftrule=3pt, arc=2pt, left=10pt, right=10pt, top=7pt, bottom=7pt}
\newtcolorbox{answerbox}{enhanced, breakable, colback=white, colframe=linec, boxrule=0.5pt, arc=2pt, left=10pt, right=10pt, top=7pt, bottom=7pt}
\newcommand{\lessonbanner}[3]{%
  \begin{tcolorbox}[enhanced, colback=indigoc, colframe=indigoc, boxrule=0pt, arc=0pt,
    left=14pt, right=14pt, top=14pt, bottom=10pt, borderline south={2.4pt}{0pt}{goldc}, fontupper=\color{white}]
    {\sffamily\footnotesize\color{white!85!black} #1}\\[4pt]{\sffamily\bfseries\Large #2}\\[3pt]{\sffamily\small\color{white!88!black} #3}
  \end{tcolorbox}}
\pagestyle{fancy}\fancyhf{}\renewcommand{\headrulewidth}{0pt}
\fancyfoot[L]{\footnotesize\color{inksoft} 공통수학2 · 2026학년도 2학기 · 지평선고등학교}
\fancyfoot[R]{\footnotesize\color{inksoft} 교사용 지도안 -- 배포 금지}

\newcommand{\mapfour}[1]{%
  \begin{tikzpicture}[scale=0.6, baseline]
    \draw[thick, indigosoft] (0,0) ellipse (0.9 and 1.6);
    \draw[thick, indigosoft] (2.6,0) ellipse (0.9 and 1.6);
    \node[circle, fill=indigoc, inner sep=1.4pt, label=left:{\footnotesize $1$}] (x1) at (0,1) {};
    \node[circle, fill=indigoc, inner sep=1.4pt, label=left:{\footnotesize $2$}] (x2) at (0,0) {};
    \node[circle, fill=indigoc, inner sep=1.4pt, label=left:{\footnotesize $3$}] (x3) at (0,-1) {};
    \node[circle, fill=goldstrong, inner sep=1.4pt, label=right:{\footnotesize $4$}] (y1) at (2.6,1) {};
    \node[circle, fill=goldstrong, inner sep=1.4pt, label=right:{\footnotesize $5$}] (y2) at (2.6,0) {};
    \node[circle, fill=goldstrong, inner sep=1.4pt, label=right:{\footnotesize $6$}] (y3) at (2.6,-1) {};
    \node[above] at (0,1.75) {\footnotesize $X$};
    \node[above] at (2.6,1.75) {\footnotesize $Y$};
    #1
  \end{tikzpicture}}

\begin{document}
\lessonbanner{공통수학2 · Ⅲ. 함수와 그래프 · 1. 함수}{01 함수 --- 1/2차시}{함수의 뜻 · 교사용 지도안 (2026학년도 2학기)}
\vspace{10pt}

\noindent\begin{tabularx}{\linewidth}{@{}>{\bfseries\color{indigoc}}p{2.6cm}X@{}}
\toprule
대상 & 고1 · 2개 반 · 반당 13명 \\
차시 & 50분 (도입 10 · 전개 30 · 정리 10) \\
성취기준 & 함수의 개념을 이해하고, 함수의 정의역, 공역, 치역을 구별할 수 있다. \\
선행 학습 & 대단원 Ⅱ 전체(집합과 명제) --- 특히 집합의 원소와 대응 관계 \\
\bottomrule
\end{tabularx}

\section{학습 목표 \& Big Idea}
\begin{ideabox}
\textbf{\color{goldstrong}영속적 이해 (Big Idea)}\\
함수란 정의역의 모든 원소 하나하나에 공역의 원소가 \textbf{빠짐없이, 그리고 오직 하나씩} 대응하는 특별한 대응 규칙이다. 이 `존재성'과 `유일성' 두 조건이 동시에 성립할 때만 대응은 함수가 된다.
\vspace{4pt}
\small\color{inksoft} 핵심개념: \textbf{대응의 유일성} \quad\textbullet\quad 핵심개념: \textbf{정의역·공역·치역} \quad\textbullet\quad 일반화: \textbf{함수는 모든 원소에 오직 하나씩만 대응하는 특별한 대응이다}
\end{ideabox}
\begin{calloutbox}
\textbf{차시 학습 목표}
\begin{itemize}[leftmargin=1.4em, itemsep=2pt, topsep=4pt]
  \item 두 집합 사이의 대응 중에서 함수인 것과 함수가 아닌 것을 구별하고 그 이유를 설명할 수 있다.
  \item 함수의 정의역, 공역, 치역을 구별하여 말하고, 함숫값을 구할 수 있다.
\end{itemize}
\end{calloutbox}

\section{질문의 연결}
\begin{tabularx}{\linewidth}{@{}>{\bfseries\color{indigoc}\small}p{2.4cm}X@{}}
사실적 질문 & 두 집합 $X$, $Y$ 사이의 대응이 함수가 되려면 어떤 조건을 만족해야 할까? \\[6pt]
개념적 질문 & 공역의 원소가 정의역의 어떤 원소와도 대응되지 않고 남아 있어도 왜 함수라고 할 수 있을까? \\[6pt]
논쟁적 질문 & ``우리 반 학생 각자에게 가장 친한 친구 한 명을 대응시키자''는 규칙은 함수라고 할 수 있을까? \\
\end{tabularx}

\section{수업 흐름}
\begin{tabularx}{\linewidth}{@{}>{\centering\arraybackslash}p{1.6cm}X>{\raggedright\arraybackslash\small}p{3.1cm}@{}}
\toprule
\textbf{단계} & \textbf{활동 \& 발문} & \textbf{ATL / VTR} \\
\midrule
\textcolor{goldstrong}{\textbf{도입}}\newline\footnotesize 10분 &
\textbf{마음공부 (2분)} --- 새 대단원 시작, 유무념 대조.\newline
\textbf{생각 열기 (8분)} --- 학생 각자에게 사물함 번호를 하나씩 대응시키는 규칙과, 학생 각자에게 좋아하는 급식 메뉴를 대응시키는 규칙(어떤 학생은 여러 개를 고르고 어떤 학생은 하나도 못 고름)을 비교 제시. 두 대응이 무엇이 다른지 토의. &
ATL: 사고(비교)\newline VTR: See-Think-Wonder \\
\addlinespace
\textcolor{goldstrong}{\textbf{전개}}\newline\footnotesize 30분 &
\textbf{대응과 함수의 정의 (10분)} --- 대응 정의 $\to$ 활동지의 네 가지 대응 그림 제시 $\to$ 문제 1(함수인 것 판별, 존재성·유일성 조건으로 이유 설명).\newline
\textbf{정의역·공역·치역 (12분)} --- 정의 도입($X$: 정의역, $Y$: 공역, $f(X)=\{f(x)\mid x\in X\}$: 치역) $\to$ 예제($X=\{1,2,3,4\}$, $Y=\{2,4,6,8,10\}$, $f(x)=2x$의 정의역·공역·치역·함숫값) $\to$ 문제 2.\newline
\textbf{서로 같은 함수 (8분)} --- 두 함수가 같으려면 정의역과 공역이 각각 같고, 정의역의 모든 원소에서 함숫값이 같아야 함을 소개 $\to$ 문제 3. &
ATT: 촉진자 역할\newline ATL: 대수적 조작·논리적 판별 \\
\addlinespace
\textcolor{goldstrong}{\textbf{정리}}\newline\footnotesize 10분 &
\textbf{개념 연결 질문 (5분)} --- ``정의역의 한 원소에 공역의 원소 두 개가 대응되면 왜 함수가 아닐까?''\newline
\textbf{성찰 (5분)} --- 감사 일기 1문장 + 다음 차시 예고(함수의 그래프, 여러 가지 함수). &
마음공부 · 감사 일기 \\
\bottomrule
\end{tabularx}

\begin{calloutbox}
\textbf{지도상의 유의점} --- 함수가 되기 위한 두 조건, 즉 정의역의 모든 원소가 빠짐없이 대응되어야 한다는 \textbf{존재성}과 각 원소에 오직 하나의 상만 대응되어야 한다는 \textbf{유일성}을 명확히 구분해 지도한다. 치역은 항상 공역의 부분집합이며, 공역과 치역이 반드시 같을 필요는 없음을 강조한다.
\end{calloutbox}

\section{형성평가 \& 답안}
\begin{answerbox}
\textbf{문제 1} \quad 다음 네 가지 대응 중 함수인 것을 모두 고르고 이유를 설명하시오.

\vspace{6pt}
\noindent
\begin{minipage}{0.24\linewidth}\centering
\mapfour{
  \draw[->, thick] (x1) -- (y1);
  \draw[->, thick] (x2) -- (y2);
  \draw[->, thick] (x3) -- (y2);
}\\ \footnotesize ①
\end{minipage}\hfill
\begin{minipage}{0.24\linewidth}\centering
\mapfour{
  \draw[->, thick] (x1) -- (y1);
  \draw[->, thick] (x3) -- (y3);
}\\ \footnotesize ②
\end{minipage}\hfill
\begin{minipage}{0.24\linewidth}\centering
\mapfour{
  \draw[->, thick] (x1) -- (y1);
  \draw[->, thick] (x1) -- (y2);
  \draw[->, thick] (x2) -- (y2);
  \draw[->, thick] (x3) -- (y3);
}\\ \footnotesize ③
\end{minipage}\hfill
\begin{minipage}{0.24\linewidth}\centering
\mapfour{
  \draw[->, thick] (x1) -- (y1);
  \draw[->, thick] (x2) -- (y1);
  \draw[->, thick] (x3) -- (y3);
}\\ \footnotesize ④
\end{minipage}

\vspace{6pt}
\textbf{\color{goldstrong}답} \; 함수: ①, ④ ($X$의 모든 원소가 빠짐없이 $Y$의 원소 하나에만 대응). \; ②는 함수 아님($2$가 어디에도 대응되지 않음, 존재성 위반). \; ③은 함수 아님($1$이 $4$와 $5$ 두 곳에 대응, 유일성 위반).
\end{answerbox}

\begin{minipage}[t]{0.48\linewidth}
\begin{answerbox}
\textbf{문제 2} \quad $X=\{1,2,3,4\}$, $Y=\{2,4,6,8,10\}$, $f(x)=2x$\\[4pt]
\textbf{\color{goldstrong}답} \; 정의역 $=\{1,2,3,4\}$, 공역 $=\{2,4,6,8,10\}$, 치역 $=\{2,4,6,8\}$, $f(3)=6$
\end{answerbox}
\end{minipage}\hfill
\begin{minipage}[t]{0.48\linewidth}
\begin{answerbox}
\textbf{문제 3} \quad $X=\{-1,0,1\}$에서 정의된 $f(x)=x^2$, $g(x)=|x|$가 서로 같은 함수인지 판별\\[4pt]
\textbf{\color{goldstrong}답} \; $f(-1)=1=g(-1)$, $f(0)=0=g(0)$, $f(1)=1=g(1)$이고 정의역·공역이 같으므로 $f=g$ (같은 함수).
\end{answerbox}
\end{minipage}

\section{성취수준 (이 차시 범위)}
\begin{tabularx}{\linewidth}{@{}>{\bfseries\color{goldstrong}}p{0.9cm}X@{}}
\toprule
A & 함수의 뜻(존재성·유일성)을 근거를 들어 설명하고, 정의역·공역·치역을 정확히 구별하며 다양한 대응에서 함수 여부를 판별할 수 있다. \\
B & 함수의 뜻을 이해하고, 정의역·공역·치역을 구하고 함숫값을 계산할 수 있다. \\
C & 함수의 뜻을 알고, 안내에 따라 정의역·공역·치역을 구할 수 있다. \\
D & 함수인 대응과 함수가 아닌 대응을 구별할 수 있다. \\
E & 두 집합 사이의 대응이 무엇인지 안다. \\
\bottomrule
\end{tabularx}

\section{다음 차시}
\begin{calloutbox}
\textbf{01 함수 --- 2/2차시.} 함수의 그래프와 여러 가지 함수(일대일함수, 일대일대응, 항등함수, 상수함수)를 다룬다.
\end{calloutbox}
\end{document}
